\documentclass[11pt,a4paper]{article}
\usepackage[margin=0.9in]{geometry}
\usepackage[T1]{fontenc}\usepackage{lmodern}\usepackage[hidelinks]{hyperref}\usepackage{microtype}
\pagestyle{empty}\setlength{\parindent}{0pt}\setlength{\parskip}{9pt}
\begin{document}
Dear Editors,

Please consider my manuscript, ``Discriminating mechanisms of ultrasensitivity in bacterial nitrogen signaling,'' as a Research Article in PLOS Computational Biology.

Similar biochemical dose--response curves can support different molecular explanations. This study uses the GlnD--PII--GlnE--GS cascade to connect explicit mechanistic ambiguity to quantitative limits on interpretation and to the choice of the next measurement.

The manuscript develops three linked results. First, finite GlnD and GlnE binding--catalysis transformations preserve complete stationary target-state distributions or downstream responses. The ambiguity therefore persists beyond fitting the same mean titration. Second, a specified protein-conserving cascade with all PII and GS modification-count states bounds differences between total observations of free-state-equivalent mechanisms. Inverting these bounds gives sufficient enzyme-to-target conditions for controlling interpretation error. Third, the unresolved mechanism determines the informative assay: state-resolved PII addresses ladder and population alternatives, while target-free GlnD occupancy and normalized GlnE binding distinguish the specified regulatory families after nuisance variation.

The contribution is the connection between these biochemical constructions, their reaction and observation boundaries, and conditional measurement design. Response decomposition and general identifiability principles are treated as established foundations. Productive regulatory complexes, conserved rather than free glutamine, and hidden conformations delimit the conclusions. The analysis does not claim to identify an unrestricted native mechanism or validate the complete physiological nitrogen network.

The work combines exact derivations, reanalysis of published measurements, and independently checked numerical reaction models. No new experimental measurements were generated. Digitized coordinates, prediction grids, parameter provenance, and verification outputs are supplied as S1 Data; reproducible analysis code is publicly available through GitHub and its versioned Zenodo archive.

I believe the explicit link from model ambiguity to experimentally testable observation choices will interest computational biologists studying signaling, parameter identifiability, and experimental design.

Thank you for considering this manuscript.

Sincerely,\\Yule Choi\\Korea Science Academy of KAIST\\Busan, Republic of Korea\\25-117@ksa.hs.kr
\end{document}
